\section{Conclusion}

The experiments in this paper support a simple claim. Neural networks maintain homeostatic equilibrium during learning and they do so through active variance redistribution across components. When I perturb a suppressor head other heads compensate. When I block that compensation learning is always harmed, either by slowing dramatically or by converging to brittle non generalizing solutions. This is Le Chatelier's principle in silicon. The system resists perturbation and works to restore a preferred equilibrium, and any intervention has to be understood in that light.
